\documentclass{dlproblemset}
\usepackage{booktabs}
\usepackage{hyperref}
\usepackage{tikz}
\usetikzlibrary{arrows.meta,decorations.pathreplacing,positioning}
\usepackage{pgfplots}
\pgfplotsset{compat=1.18}
\usepackage{emoji}
\tikzset{every picture/.style={line width=0.75pt}}

\definecolor{CellBlue}{HTML}{BBDEFB}
\definecolor{CellGreen}{HTML}{C8E6C9}
\definecolor{CellYellow}{HTML}{FFF9C4}
\definecolor{CellRed}{HTML}{FFCDD2}
\definecolor{DarkBlue}{HTML}{1565C0}
\definecolor{DarkGreen}{HTML}{1B5E20}
\definecolor{DarkRed}{HTML}{C62828}

\title{Atenção e Transformers}
\subtitle{Decoder Only Models}

\begin{document}

\paragraph{Instruções:} Cada questão possui exatamente uma alternativa correta.


\section*{Questões}
% ============================================================

\begin{enumerate}[I)]

% ------------------------------------------------------------------
% Q1 — Fácil (KV-Cache)
% ------------------------------------------------------------------
\item Durante a fase de \textbf{decode} de um decoder-only com KV-cache,
      o modelo gera o quarto token de uma sequência iniciada com o
      prompt de 3 tokens \textit{``O rato roeu''}. O diagrama abaixo
      mostra o estado dos pares K,V nesse passo.

\begin{center}
\begin{tikzpicture}[
  tok/.style={draw, rounded corners, minimum width=1.3cm, minimum height=0.46cm,
              align=center, font=\small},
  kv/.style={draw, minimum width=1.3cm, minimum height=0.36cm,
             align=center, font=\scriptsize},
  arr/.style={->, >=Stealth, thick},
]
  \node[tok, fill=CellBlue]   (t1) at (0.0, 2.5) {\texttt{O}};
  \node[tok, fill=CellBlue]   (t2) at (1.7, 2.5) {\texttt{rato}};
  \node[tok, fill=CellBlue]   (t3) at (3.4, 2.5) {\texttt{roeu}};
  \node[tok, fill=CellRed]    (t4) at (5.1, 2.5) {\texttt{?}};

  \node[kv, fill=CellGreen]   (k1) at (0.0, 1.58) {$K_1$};
  \node[kv, fill=CellGreen]   (k2) at (1.7, 1.58) {$K_2$};
  \node[kv, fill=CellGreen]   (k3) at (3.4, 1.58) {$K_3$};
  \node[kv, fill=CellYellow]  (k4) at (5.1, 1.58) {$K_4$};

  \node[kv, fill=CellGreen]   (v1) at (0.0, 0.96) {$V_1$};
  \node[kv, fill=CellGreen]   (v2) at (1.7, 0.96) {$V_2$};
  \node[kv, fill=CellGreen]   (v3) at (3.4, 0.96) {$V_3$};
  \node[kv, fill=CellYellow]  (v4) at (5.1, 0.96) {$V_4$};

  \node[font=\scriptsize, gray, left=0.1cm of k1] {$K$:};
  \node[font=\scriptsize, gray, left=0.1cm of v1] {$V$:};

  \draw[arr] (t4.south) -- (k4.north);

  \node[font=\scriptsize, below=0.1cm] at (2.55, 0.6)
    {\colorbox{CellGreen}{\,}~cache~~\colorbox{CellYellow}{\,}~computando agora};
\end{tikzpicture}
\end{center}

O que o diagrama ilustra sobre o funcionamento do KV-cache nesse passo?

\begin{enumerate}[a)]
    \item Apenas $K_4$ e $V_4$ são computados; $K_1$--$K_3$ e $V_1$--$V_3$
          são reutilizados do cache.
    \item Todos os $K_1$--$K_4$ e $V_1$--$V_4$ são recalculados integralmente
          a cada novo token gerado.
    \item $K_4$ é computado, mas todos os $V_1$--$V_4$ são recalculados
          para combinar com a nova query $Q_4$.
    \item O cache armazena apenas as queries; keys e values são sempre
          recalculados a partir dos tokens originais.
\end{enumerate}

% ------------------------------------------------------------------
% Q2 — Fácil (Temperatura e estratégias de amostragem)
% ------------------------------------------------------------------
\item O gráfico abaixo mostra a distribuição de probabilidade sobre
      quatro tokens para os mesmos logits $z{=}[3,\;1,\;0,\;{-1}]$,
      com três valores de temperatura.

\begin{center}
\begin{tikzpicture}
  \begin{axis}[
    ybar, bar width=0.22cm,
    xlabel={Token}, ylabel={Probabilidade},
    symbolic x coords={$v_1$,$v_2$,$v_3$,$v_4$},
    xtick=data,
    ymin=0, ymax=1.08,
    width=8.5cm, height=4.8cm,
    enlarge x limits=0.20,
    legend style={at={(0.98,0.97)}, anchor=north east, font=\small},
    legend cell align=left,
  ]
    \addplot[fill=CellBlue,   draw=DarkBlue]  coordinates
      {($v_1$,0.979)($v_2$,0.018)($v_3$,0.002)($v_4$,0.000)};
    \addlegendentry{$T{=}0.5$}
    \addplot[fill=CellYellow, draw=DarkRed]   coordinates
      {($v_1$,0.831)($v_2$,0.112)($v_3$,0.041)($v_4$,0.015)};
    \addlegendentry{$T{=}1$}
    \addplot[fill=CellGreen,  draw=DarkGreen] coordinates
      {($v_1$,0.579)($v_2$,0.213)($v_3$,0.129)($v_4$,0.079)};
    \addlegendentry{$T{=}2$}
  \end{axis}
\end{tikzpicture}
\end{center}

Qual afirmação descreve corretamente o efeito da temperatura?

\begin{enumerate}[a)]
    \item $T < 1$ achata a distribuição, aumentando a diversidade;
          $T > 1$ concentra a massa no token mais provável.
    \item $T = 1$ produz distribuição uniforme; valores maiores
          concentram a massa progressivamente.
    \item $T \to 0$ equivale ao argmax (greedy); $T > 1$ achata a
          distribuição, aumentando a diversidade da geração.
    \item A temperatura não altera as razões de probabilidade entre
          tokens: se $P(v_i)/P(v_j) = r$ com $T{=}1$, essa razão
          permanece $r$ para qualquer valor de $T$.
\end{enumerate}

% ------------------------------------------------------------------
% Q3 — Médio (Perplexidade vs. benchmarks de avaliação)
% ------------------------------------------------------------------
\item Ao avaliar um LLM, um pesquisador considera usar a
      \textbf{perplexidade} sobre um corpus de teste ou o benchmark
      \textbf{MMLU} (conhecimento enciclopédico em 57 domínios,
      formato de múltipla escolha). Qual afirmação descreve
      corretamente uma vantagem e uma limitação de cada métrica?

\begin{enumerate}[a)]
    \item A perplexidade requer um conjunto de questões rotuladas,
          tornando-a mais custosa que o MMLU; o MMLU, por ser gerado
          automaticamente, é imune a contaminação do pré-treinamento.
    \item A perplexidade independe do tokenizador, permitindo
          comparação direta entre modelos com vocabulários distintos;
          o MMLU mede exclusivamente raciocínio matemático em
          múltiplos passos.
    \item A perplexidade é barata de computar e não exige dados
          rotulados, mas correlaciona-se imperfeitamente com o
          desempenho em tarefas específicas; o MMLU mede capacidades
          concretas, mas pode ser inflacionado por contaminação do
          corpus de pré-treinamento.
    \item A perplexidade garante que modelos com valores menores
          terão melhor desempenho em qualquer tarefa downstream;
          benchmarks como o MMLU são imunes a esse problema por
          usarem questões de múltipla escolha.
\end{enumerate}

% ------------------------------------------------------------------
% Q4 — Médio (MHA / MQA / GQA: tamanho de KV-cache)
% ------------------------------------------------------------------
\item Considere um modelo com $H = 8$ cabeças de query, $d_k = 64$
      e $L = 32$ camadas. O diagrama abaixo compara MHA e GQA
      com $G = 2$ grupos (usando $H{=}4$ na visualização):

\begin{center}
\begin{tikzpicture}[
  qh/.style={draw, fill=CellBlue, rounded corners,
             minimum width=0.72cm, minimum height=0.58cm,
             font=\scriptsize, align=center},
  kh/.style={draw, fill=CellGreen, rounded corners,
             minimum width=0.80cm, minimum height=0.58cm,
             font=\scriptsize, align=center},
  arr/.style={->, >=Stealth, semithick},
]
  %--- MHA ---
  \node[font=\small\bfseries] at (1.35, 2.7) {MHA ($G{=}H{=}4$)};
  \foreach \i/\x in {1/0.3, 2/1.0, 3/1.7, 4/2.4} {
    \node[qh] (qm\i) at (\x, 1.8) {Q\textsubscript{\i}};
    \node[kh] (km\i) at (\x, 0.8) {KV\textsubscript{\i}};
    \draw[arr] (km\i) -- (qm\i);
  }
  \node[font=\scriptsize, gray] at (1.35, 0.18) {4 pares KV};

  %--- GQA (G=2) ---
  \node[font=\small\bfseries] at (5.3, 2.7) {GQA ($G{=}2$)};
  \foreach \i/\x in {1/4.1, 2/4.9, 3/5.7, 4/6.5} {
    \node[qh] (qg\i) at (\x, 1.8) {Q\textsubscript{\i}};
  }
  \node[kh] (kga) at (4.5, 0.8) {KV\textsubscript{a}};
  \node[kh] (kgb) at (6.1, 0.8) {KV\textsubscript{b}};
  \draw[arr] (kga) -- (qg1);
  \draw[arr] (kga) -- (qg2);
  \draw[arr] (kgb) -- (qg3);
  \draw[arr] (kgb) -- (qg4);
  \node[font=\scriptsize, gray] at (5.3, 0.18) {2 pares KV};
\end{tikzpicture}
\end{center}

Usando $H{=}8$, $G{=}2$, $d_k{=}64$, $L{=}32$: qual é o tamanho do
KV-cache \textbf{por token} com GQA e quantas vezes menor é que o MHA?

\begin{enumerate}[a)]
    \item $8\,192$ floats/token; $4\times$ menor que MHA.
    \item $8\,192$ floats/token; $8\times$ menor que MHA.
    \item $16\,384$ floats/token; $2\times$ menor que MHA.
    \item $4\,096$ floats/token; $8\times$ menor que MHA.
\end{enumerate}

% ------------------------------------------------------------------
% Q5 — Médio (RMSNorm vs. LayerNorm)
% ------------------------------------------------------------------
\item Dado $\mathbf{x} = [3,\; 4]$ e $\boldsymbol{\gamma} = [1,\; 1]$
      (ignorando $\epsilon$), calcule $\mathrm{RMSNorm}(\mathbf{x})$
      e identifique qual das opções corresponde ao resultado correto.
      \[
        \mathrm{RMSNorm}(\mathbf{x}) = \frac{\mathbf{x}}{\mathrm{RMS}(\mathbf{x})} \odot \boldsymbol{\gamma},
        \qquad
        \mathrm{RMS}(\mathbf{x}) = \sqrt{\frac{1}{d}\sum_{i=1}^d x_i^2}.
      \]

\begin{enumerate}[a)]
    \item $[\,-1.00,\;+1.00\,]$
    \item $[\,0.849,\;1.131\,]$
    \item $[\,0.600,\;0.800\,]$
    \item $[\,0.857,\;1.143\,]$
\end{enumerate}

% ------------------------------------------------------------------
% Q6 — Médio (MoE: Expert Collapse)
% ------------------------------------------------------------------
\item Em um modelo Mixture-of-Experts (MoE) com $E = 4$ especialistas
      e $k = 1$ (top-1 routing), observa-se que, após o treinamento,
      o especialista~1 processa ${\approx}\,97\%$ de todos os tokens,
      enquanto os especialistas 2--4 ficam essencialmente inativos.
      Qual fenômeno descreve essa situação e qual é a solução adotada?

\begin{enumerate}[a)]
    \item \textbf{Overfitting} do router: o router memoriza os tokens de treino;
          solução: aumentar o número de especialistas $E$.
    \item \textbf{Expert collapse}: ciclo de reforço positivo em que o especialista
          mais ativo recebe mais tokens, melhora mais e atrai ainda mais tokens;
          solução: perda auxiliar de balanceamento
          $\mathcal{L}_{\text{aux}} = \alpha\,E\sum_{i=1}^{E} f_i\,P_i$.
    \item \textbf{Vanishing gradients}: gradientes desaparecem nos especialistas
          inativos; solução: conexões residuais no router.
    \item \textbf{Mode collapse}: todos os especialistas convergem para a mesma
          função; solução: inicialização diversificada dos pesos.
\end{enumerate}

% ------------------------------------------------------------------
% Q7 — Médio (SwiGLU: mecanismo de portão)
% ------------------------------------------------------------------
\item Na FFN SwiGLU, o cálculo é:
      \[
        \mathrm{FFN}(x) = \bigl(\mathrm{SiLU}(xW_1) \odot xW_3\bigr)W_2,
      \]
      enquanto a ReLU-FFN clássica usa $\mathrm{FFN}(x) = \max(0,\,xW_1)\,W_2$.
      Qual é a principal vantagem do portão $xW_3$ em relação ao corte
      estático da ReLU?

\begin{enumerate}[a)]
    \item A SwiGLU usa três matrizes em vez de duas, aumentando a capacidade
          bruta sem custo computacional adicional por token.
    \item A SwiGLU aplica a não-linearidade duas vezes consecutivas
          ($\mathrm{SiLU}$ e depois $W_2$), aumentando a profundidade efetiva.
    \item O portão $xW_3$ aprende dinamicamente quais dimensões do espaço
          oculto deixar passar, ao contrário da ReLU, que zera estaticamente
          valores negativos independentemente do contexto da entrada.
    \item A SwiGLU elimina a necessidade de normalização por camada, pois
          o portão já controla a magnitude das ativações.
\end{enumerate}

% ------------------------------------------------------------------
% Q8 — �� Difícil (Decodificação Especulativa)
% ------------------------------------------------------------------
\item \emoji{skull-and-crossbones} No algoritmo de decodificação especulativa
      (Leviathan et al., 2023), o modelo alvo $M_p$ e o modelo draft $M_q$
      atribuem as seguintes probabilidades a três tokens candidatos:

\begin{center}
\begin{tabular}{lccc}
  \toprule
  Token & $p$ (alvo $M_p$) & $q$ (draft $M_q$) \\
  \midrule
  $\tilde{x}$ & 0.30 & 0.60 \\
  $\tilde{y}$ & 0.50 & 0.30 \\
  $\tilde{z}$ & 0.20 & 0.10 \\
  \bottomrule
\end{tabular}
\end{center}

O draft propõe $\tilde{x}$. Qual é a probabilidade de aceitação de
$\tilde{x}$ e, caso rejeitado, de qual distribuição é re-amostrado o
token substituto?

\begin{enumerate}[a)]
    \item Aceitação $= 1.00$ (sempre aceito, pois $p(\tilde{x}) > 0$);
          substituto amostrado de $p$ sem modificação.
    \item Aceitação $= 0.50$; substituto amostrado de
          $p'(x) \propto \max\!\bigl(0,\,p(x)-q(x)\bigr)$,
          que concentra a massa onde $M_p$ supera $M_q$.
    \item Aceitação $= 0.50$; substituto amostrado
          uniformemente sobre o vocabulário.
    \item Aceitação $= 0.30$; substituto amostrado de $q$ (draft).
\end{enumerate}

% ------------------------------------------------------------------
% Q9 — Difícil (Leis de Escala / Chinchilla: regra 20N e retornos decrescentes)
% ------------------------------------------------------------------
\item Uma equipe treina um modelo com $N = 8 \times 10^9$ parâmetros durante
      $D = 80 \times 10^9$ tokens e obtém perda de validação de $2.22$.
      Satisfeitos com o resultado, param o treinamento.
      Em seguida, avaliam se vale a pena uma segunda rodada de treinamento
      prolongado até $D = 800 \times 10^9$ tokens.
      À luz das leis de escala do Chinchilla (Hoffmann et al., 2022),
      qual afirmação descreve corretamente a situação?

\begin{enumerate}[a)]
    \item O modelo foi treinado de forma ótima ($D = 10N$ é a
          regra do Chinchilla); a segunda rodada não trará nenhuma melhoria
          pois a perda já saturou em $2.22$.
    \item O modelo foi sub-treinado em dados ($D^* = 20N = 160\text{B} > D$);
          a segunda rodada ainda reduzirá a perda, mas com
          \emph{retornos decrescentes}.
    \item O modelo foi super-treinado em dados ($D = 80\text{B} > D^* = N = 8\text{B}$);
          a segunda rodada causará sobreajuste e a perda de validação
          aumentará.
    \item O modelo foi treinado de forma ótima ($D = 10N$);
          a segunda rodada melhorará a perda linearmente com os tokens
          adicionados.
\end{enumerate}

% ------------------------------------------------------------------
% Q10 — Médio (Few-shot prompting: restrição de janela de contexto)
% ------------------------------------------------------------------
\item Uma pesquisadora usa $k$-shot prompting com um LLM de janela de
      contexto $C = 2048$ tokens. O prompt de sistema ocupa $150$ tokens,
      cada exemplar (pergunta $+$ rótulo) ocupa $200$ tokens e a consulta
      de teste ocupa $200$ tokens. Com $k = 5$ o prompt cabe na janela.
      Ela quer testar $k = 10$. Qual afirmação está correta?

\begin{enumerate}[a)]
    \item $k = 10$ é viável: $10 \times 200 = 2000 < 2048$, portanto
          há tokens suficientes.
    \item $k = 10$ não é viável: o prompt totalizaria
          $150 + 10 \times 200 + 200 = 2350 > 2048$ tokens, excedendo a
          janela; o modelo precisaria truncar o contexto ou a geração.
    \item $k = 10$ é equivalente a \textit{fine-tuning} com 10 exemplos:
          em ambos os casos, o modelo atualiza sua distribuição de saída
          sem modificar os pesos.
    \item $k = 10$ é viável, mas degrada o desempenho: mais exemplares no
          contexto causam \textit{overfitting} aos rótulos, análogo ao
          fine-tuning com poucos dados.
\end{enumerate}
% R: b

% ------------------------------------------------------------------
% Q11 — Médio (MQA: redução de KV-cache)
% ------------------------------------------------------------------
\item O \textbf{Multi-Query Attention} (MQA, Shazeer 2019) substitui os
      $H$ pares $K,V$ independentes do MHA por um \textbf{único} par $K,V$
      compartilhado entre todas as cabeças de query.
      Considere um modelo com $H = 8$ cabeças, $d_k = 64$ e $L = 32$ camadas.
      Qual é o tamanho do KV-cache \textbf{por token} com MQA e por quantas
      vezes é menor que o MHA equivalente?

\begin{enumerate}[a)]
    \item $16\,384$ floats/token; $2\times$ menor que MHA.
    \item $4\,096$ floats/token; $8\times$ menor que MHA.
    \item $8\,192$ floats/token; $4\times$ menor que MHA
          (mesmo resultado do GQA com $G{=}2$).
    \item $4\,096$ floats/token; $4\times$ menor que MHA.
\end{enumerate}

% ------------------------------------------------------------------
% Q12 — Médio (Top-p / Nucleus Sampling)
% ------------------------------------------------------------------
\item Num passo de decodificação, o modelo atribui as seguintes
      probabilidades (ordenadas de forma decrescente):

\begin{center}
\begin{tabular}{lcc}
  \toprule
  Token & $P$ & $P$ acumulada \\
  \midrule
  \texttt{o}    & 0.45 & 0.45 \\
  \texttt{um}   & 0.30 & 0.75 \\
  \texttt{este} & 0.15 & 0.90 \\
  \texttt{seu}  & 0.07 & 0.97 \\
  \texttt{...}  & 0.03 & 1.00 \\
  \bottomrule
\end{tabular}
\end{center}

Com \textbf{top-p = 0.90} (nucleus sampling), qual é o conjunto
\emph{mínimo} de tokens que compõe o núcleo do qual o próximo token
será amostrado?

\begin{enumerate}[a)]
    \item Apenas \texttt{o} (probabilidade acumulada = 0.45).
    \item \texttt{o}, \texttt{um}
          (probabilidade acumulada = 0.75 $<$ 0.90).
    \item \texttt{o}, \texttt{um}, \texttt{este}
          (probabilidade acumulada = 0.90 $\ge$ 0.90).
    \item \texttt{o}, \texttt{um}, \texttt{este}, \texttt{seu}
          (probabilidade acumulada = 0.97 $>$ 0.90).
\end{enumerate}

% ------------------------------------------------------------------
% Q13 — Médio (Perplexidade)
% ------------------------------------------------------------------
\item Um modelo de linguagem com vocabulário de 6 tokens
      $\mathcal{V} = \{\texttt{o},\;\texttt{gato},\;\texttt{correu},\;
      \texttt{leite},\;\texttt{bebeu},\;\texttt{fim}\}$
      atribui as seguintes probabilidades condicionais para cada posição
      da sequência \textit{``o~gato~correu''}:

\begin{center}
\begin{tabular}{lrrr}
  \toprule
  \textbf{Token}
    & $P(w_1)$
    & $P(w_2\mid w_1)$
    & $P(w_3\mid w_{1:2})$ \\
  \midrule
  \texttt{o}       & \textbf{0.500} & 0.050 & 0.050 \\
  \texttt{gato}    & 0.150 & \textbf{0.250} & 0.100 \\
  \texttt{correu}  & 0.150 & 0.300 & \textbf{0.125} \\
  \texttt{leite}   & 0.100 & 0.150 & 0.375 \\
  \texttt{bebeu}   & 0.050 & 0.150 & 0.250 \\
  \texttt{fim}     & 0.050 & 0.100 & 0.100 \\
  \bottomrule
\end{tabular}
\end{center}

Os valores em \textbf{negrito} correspondem ao token observado em
cada posição. Usando a definição
\[
  \mathrm{PPL}(w_1,\ldots,w_N)
    = \exp\!\left(-\frac{1}{N}\sum_{t=1}^{N}\ln P(w_t\mid w_{<t})\right),
\]
qual é a \textbf{perplexidade} do modelo nessa sequência?

\begin{enumerate}[a)]
    \item $\approx 3.4$
    \item $4$
    \item $8$
    \item $64$
\end{enumerate}

\end{enumerate}



% ============================================================
\section*{Soluções}
% ============================================================

% --- Solução Q1 ---
\subsection*{I — Resposta: (a)}
Durante a fase de \emph{decode}, o KV-cache armazena os pares $K_j, V_j$ de todos
os tokens já processados. Para gerar o token na posição~4, o modelo computa $K_4$
e $V_4$ aplicando as projeções $W^K$ e $W^V$ ao embedding do novo token; os pares
$K_1,V_1$, $K_2,V_2$ e $K_3,V_3$ são lidos diretamente do cache sem recálculo.
Isso reduz o custo por passo de decodificação de $\mathcal{O}(n^2 d)$ (sem cache)
para $\mathcal{O}(n d)$ (com cache), sendo essencial para geração eficiente em
sequências longas.
% ---

% --- Solução Q2 ---
\subsection*{II — Resposta: (c)}
A temperatura escala os logits antes do softmax: $P_T(v) \propto \exp(z_v/T)$.
Com $T \to 0$, apenas o maior logit contribui significativamente, equivalente ao
$\arg\max$ (greedy decoding). Com $T > 1$, os logits são reduzidos em magnitude,
tornando a distribuição mais uniforme e aumentando a diversidade das amostras.
O gráfico confirma: $T{=}0.5$ (azul) concentra ${\approx}97.9\%$ em $v_1$;
$T{=}2$ (verde) distribui $57.9\%$ em $v_1$ e o restante nos demais tokens.
A opção~(a) inverte os efeitos. A opção~(b) confunde $T{=}0$ com distribuição
uniforme. A opção~(d) é incorreta pois a temperatura altera as probabilidades
relativas: a razão $P_T(v_i)/P_T(v_j) = e^{(z_i-z_j)/T}$ depende de $T$.
% ---

% --- Solução Q3 ---
\subsection*{III — Resposta: (c)}
A perplexidade ($\mathrm{PP} = \exp\!\bigl(-\tfrac{1}{N}\sum_{i=1}^{N}\log P(w_i)\bigr)$)
não exige anotação humana e é eficiente de calcular, sendo amplamente usada como
proxy de qualidade durante o desenvolvimento e treinamento. Sua limitação principal
é a correlação imperfeita com tarefas downstream: modelos com perplexidade similar
podem diferir substancialmente em benchmarks como MMLU ou GSM8k. O MMLU
(Hendrycks et al., 2021) avalia conhecimento em 57 domínios via múltipla escolha,
fornecendo uma medida interpretável de capacidades específicas; porém, por ser
público e amplamente divulgado, suas questões podem aparecer no corpus de
pré-treinamento (\emph{data contamination}), inflando artificialmente o desempenho
reportado. A opção~(a) inverte o requisito de dados rotulados: perplexidade não os
exige, o MMLU sim. A opção~(b) confunde dependência do tokenizador (perplexidade
\emph{depende} do tokenizador) e atribui ao MMLU o foco matemático do GSM8k.
A opção~(d) afirma falsamente que perplexidade baixa garante bom desempenho em
qualquer tarefa downstream.
% ---

% --- Solução Q4 ---
\subsection*{IV — Resposta: (a)}
O tamanho do KV-cache por token é $2 \times G \times d_k \times L$ floats
(2 tensores: K e V; $G$ grupos KV; $d_k$ dimensão por cabeça; $L$ camadas):
\[
  \text{GQA}\;(G{=}2):\quad 2 \times 2 \times 64 \times 32 = 8\,192\ \text{floats/token}.
\]
Com MHA ($G = H = 8$):
\[
  2 \times 8 \times 64 \times 32 = 32\,768\ \text{floats/token}.
\]
Razão: $32\,768 / 8\,192 = 4\times$, não $8\times$, pois $G{=}2$ reduz $H$ por
um fator $H/G = 8/2 = 4$, não por $H = 8$. A opção~(d) corresponderia ao
MQA ($G{=}1$), onde há apenas 1 par KV compartilhado.
% ---

% --- Solução Q5 ---
\subsection*{V — Resposta: (b)}
\[
  \mathrm{RMS}(\mathbf{x})
    = \sqrt{\frac{3^2 + 4^2}{2}}
    = \sqrt{\frac{25}{2}}
    = \sqrt{12.5}
    \approx 3.536.
\]
\[
  \mathrm{RMSNorm}(\mathbf{x})
    = \left[\frac{3}{3.536},\;\frac{4}{3.536}\right]
    \approx [0.849,\;1.131].
\]
A opção~(a) é o resultado do \emph{LayerNorm} ($\mu = 3.5$, $\sigma = 0.5$),
que \textbf{subtrai a média} antes de normalizar; o RMSNorm não subtrai a média,
por isso a segunda componente pode ser maior que~1. A opção~(c) seria a normalização
pela norma $\ell_2$ ($\|\mathbf{x}\|_2 = 5 \Rightarrow [3/5,\,4/5] = [0.6,\,0.8]$);
a opção~(d) divide pela média $\mu = 3.5$.
% ---

% --- Solução Q6 ---
\subsection*{VI — Resposta: (b)}
O \textbf{expert collapse} é um ciclo de reforço positivo: o especialista mais
capaz recebe mais tokens $\to$ acumula mais gradiente $\to$ torna-se ainda melhor
$\to$ atrai ainda mais tokens. Os demais especialistas recebem gradiente próximo
de zero e nunca aprendem. A solução padrão é acrescentar uma perda auxiliar de
balanceamento à loss principal:
\[
  \mathcal{L}_{\text{aux}} = \alpha\, E \sum_{i=1}^{E} f_i\, P_i,
\]
onde $f_i$ é a fração de tokens roteados ao especialista $i$ e $P_i$ é o score
médio do router para esse especialista. Minimizar $\mathcal{L}_{\text{aux}}$
incentiva distribuição uniforme entre especialistas (o Mixtral usa $\alpha=0.01$).
% ---

% --- Solução Q7 ---
\subsection*{VII — Resposta: (c)}
Na ReLU-FFN, $\max(0, z) = 0$ sempre que $z < 0$, independentemente do token
de entrada; o corte é \emph{estático} e depende apenas do valor da ativação.
Na SwiGLU, o portão $xW_3$ é um vetor computado a partir da entrada $x$: para
cada token diferente, $xW_3$ produz valores distintos, de modo que o modelo
\emph{aprende} dinamicamente quais dimensões do espaço de dimensão $d_\text{ff}$
suprimir ou amplificar. Isso confere maior expressividade sem alterar a
complexidade assintótica. A opção~(a) é parcialmente verdadeira (SwiGLU usa 3
matrizes), mas isso \emph{aumenta} o custo computacional; a equivalência de
parâmetros é mantida reduzindo $d_\text{ff}$ para $\tfrac{8}{3}\,d_\text{model}$.
% ---

% --- Solução Q8 ---
\subsection*{VIII — Resposta: (b)}
\textbf{Probabilidade de aceitação:}
\[
  \min\!\left(1,\;\frac{p(\tilde{x})}{q(\tilde{x})}\right)
  = \min\!\left(1,\;\frac{0.30}{0.60}\right) = 0.50.
\]
\textbf{Distribuição de re-amostragem} (se $\tilde{x}$ for rejeitado),
$p'(x) \propto \max\!\bigl(0,\, p(x) - q(x)\bigr)$:
\begin{center}
\begin{tabular}{lcccc}
  \toprule
  Token & $p$ & $q$ & $p - q$ & $p'$ (normalizado) \\
  \midrule
  $\tilde{x}$ & 0.30 & 0.60 & $-$0.30 & 0.00 \\
  $\tilde{y}$ & 0.50 & 0.30 & $+$0.20 & $\approx$0.667 \\
  $\tilde{z}$ & 0.20 & 0.10 & $+$0.10 & $\approx$0.333 \\
  \bottomrule
\end{tabular}
\end{center}
$p'$ concentra a massa nos tokens onde $M_p$ supera $M_q$, garantindo que a
distribuição final seja \textbf{idêntica} a amostrar diretamente de $M_p$,
preservando a qualidade do modelo alvo por construção. A opção~(a) confunde
``$p > 0$'' com ``aceitar sempre''; apenas $p \ge q$ garante aceitação com
probabilidade~1.
% ---

% --- Solução Q9 ---
\subsection*{IX — Resposta: (b)}
A regra prática do Chinchilla estabelece $D^* \approx 20N$.
Para $N = 8 \times 10^9$ parâmetros:
\[
  D^* = 20 \times 8 \times 10^9 = 160 \times 10^9\ \text{tokens}.
\]
O treinamento parou em $D = 80$B $< 160$B, portanto o modelo está
\textbf{sub-treinado em dados}.

Continuar até $D = 800$B ainda melhorará o modelo:
a perda segue uma \textbf{lei de potência} $L(D) \propto D^{-\alpha_D}$
($\alpha_D \approx 0.095$), de modo que ganhos existem além de $D^*$,
mas são cada vez menores. Essa é a motivação por trás de modelos como o Llama~3 8B,
treinado em ${\approx}15$T tokens (${\approx}1900 \times N^*$): do ponto de
vista de \emph{inferência}, vale gastar mais no treinamento para obter
um modelo menor e mais barato de implantar.

As opções~(a) e~(d) erram a regra ($10N$ em vez de $20N$).
A opção~(a) agrava o erro ao afirmar que a perda satura exatamente em $20 \times$, ela apenas \emph{diminui mais lentamente}, mas não para.
A opção~(c) inverte a direção: com $D^* = 20N = 160$B, o modelo com
$D = 80$B está \emph{abaixo}, não acima, do ótimo; além disso, LLMs
não sofrem sobreajuste por excesso de dados, tipicamente quanto mais dados melhor.
% ---

% --- Solução Q10 ---
\subsection*{X — Resposta: (b)}
Com $k = 5$: $150 + 5 \times 200 + 200 = 1350 \le 2048$ tokens (cabe).
Com $k = 10$: $150 + 10 \times 200 + 200 = 2350 > 2048$ tokens (excede a janela em
$302$ tokens; o modelo precisaria truncar exemplares iniciais ou a geração, degradando a qualidade).
A opção~(a) esquece o prompt de sistema ($150$) e a consulta ($200$), computando apenas
$k \times 200 = 2000 < 2048$ e concluindo erroneamente que cabe.
A opção~(c) confunde \textit{in-context learning} com \textit{fine-tuning}: o \textit{few-shot
prompting} não atualiza nenhum peso: o modelo usa os exemplares apenas como contexto durante
a inferência.
A opção~(d) aplica indevidamente o conceito de \textit{overfitting} ao ICL: como não há
atualização de gradiente, não há sobreajuste no sentido clássico; mais exemplares tendem a
melhorar o desempenho até a janela de contexto ser atingida.
% ---

% --- Solução Q11 ---
\subsection*{XI — Resposta: (b)}
Com MQA há apenas $G = 1$ par $K,V$, compartilhado pelas $H = 8$ cabeças de query.
O tamanho do KV-cache por token é $2 \times G \times d_k \times L$:
\[
  \text{MQA}\;(G{=}1):\quad 2 \times 1 \times 64 \times 32 = 4\,096\ \text{floats/token}.
\]
O MHA usa $G = H = 8$:
\[
  \text{MHA}\;(G{=}8):\quad 2 \times 8 \times 64 \times 32 = 32\,768\ \text{floats/token}.
\]
Razão: $32\,768\,/\,4\,096 = 8\times$. O MQA é, portanto, o caso extremo da família
GQA com $G{=}1$, oferecendo a maior redução de cache ao custo de uma pequena queda
de qualidade (comparado ao GQA com $G{=}2$, que reduz $4\times$ e preserva melhor
a qualidade, ver questão~IV).
% ---

% --- Solução Q12 ---
\subsection*{XII — Resposta: (c)}
O \textbf{nucleus sampling} (Holtzman et al., 2020) define o núcleo como o menor
conjunto de tokens cuja probabilidade acumulada atinge ou supera o limiar $p$.
Com $p = 0.90$:
\begin{itemize}
  \item Após \texttt{o}: acum.\ $= 0.45 < 0.90$: insuficiente.
  \item Após \texttt{um}: acum.\ $= 0.75 < 0.90$: insuficiente.
  \item Após \texttt{este}: acum.\ $= 0.90 \ge 0.90$: limiar atingido.
\end{itemize}
O núcleo é $\{\texttt{o},\,\texttt{um},\,\texttt{este}\}$; suas probabilidades são
re-normalizadas e o próximo token é amostrado desse conjunto. A opção~(d) inclui
\texttt{seu} desnecessariamente, pois o limiar já foi satisfeito. A opção~(b) é o
top-p com $p{=}0.75$. A diferença-chave entre top-$k$ e top-$p$ é que o nucleus
adapta o tamanho do conjunto ao \emph{formato} da distribuição: se a distribuição
for concentrada (e.g., $p_1 = 0.95$), o núcleo terá poucos tokens; se for difusa,
terá muitos, comportamento mais robusto que um $k$ fixo.
% ---

% --- Solução Q13 ---
\subsection*{XIII — Resposta: (b)}
As probabilidades dos tokens observados são
$P(\texttt{o}) = \tfrac{1}{2}$,\;
$P(\texttt{gato}\mid\texttt{o}) = \tfrac{1}{4}$ e
$P(\texttt{correu}\mid\texttt{o},\texttt{gato}) = \tfrac{1}{8}$.
Aplicando a fórmula com $N = 3$:
\[
  \mathrm{PPL}
    = \exp\!\left(-\frac{1}{3}\Bigl(
        \ln\tfrac{1}{2} + \ln\tfrac{1}{4} + \ln\tfrac{1}{8}
      \Bigr)\right)
    = \exp\!\left(\frac{1 + 2 + 3}{3}\,\ln 2\right)
    = \exp\bigl(2\ln 2\bigr)
    = 2^2 = 4.
\]
\textbf{Análise das alternativas incorretas:}
\begin{itemize}
  \item \emph{(a) $\approx 3.4$}: surge ao usar a média \emph{aritmética}
        das probabilidades $(0.5 + 0.25 + 0.125)/3 \approx 0.292$
        e tomar o seu inverso; a perplexidade equivale à média
        \emph{geométrica} dos inversos, via $\exp(-\bar{\ell})$,
        não à média aritmética.
  \item \emph{(c) $= 8$}: corresponde a $1/P(w_3) = 1/0.125 = 8$,
        considerando apenas o token de maior surpresa e ignorando as
        demais posições.
  \item \emph{(d) $= 64$}: resultado de esquecer a normalização $1/N$:
        $\exp\!\bigl(-(\ln\tfrac{1}{2}+\ln\tfrac{1}{4}+\ln\tfrac{1}{8})\bigr)
        = 2^6 = 64$.
\end{itemize}
% ---

\end{document}


